\specialchapter{术语索引}

这些参考编号分别表示章、节和编号。

绝对单（李代数）：VIII.3.2 伴随（SL(2, k) 的表示）：VIII.1.4 容许（格）：VIII.12.8 单房：IX.5.2 反埃尔米特（自同态）：IX.1.1 反不变：IX.7.4 双序（李代数的）：VIII.12.1 Borel（子代数）：VIII.3.3, VIII.3.5 典则（嘉当子代数、根系、基）：VIII.5.3 典则（sl(2, k) 的对合）：VIII.1.1 嘉当（子代数）：VII.2.1 Casimir 元：IX.7.6 中心（g-模的特征标）：VIII.6.1 中心（函数）：IX.8.3 房：IX.5.2 特征标（g-模的）：VIII.7.7 Chevalley（序）：VIII.12.7 Chevalley（系）：VIII.2.4 经典（可分裂李代数）：VIII.3.2 净（子群）：IX.4, 习题 15 Clebsch-Gordan（公式）：VIII.9.4 紧（李代数）：IX.1.3 紧（复李代数的实形式）：IX.3.1 相容（表示）：VII.3.1, VIII.1.4 条件 (AC)：VII.1.1 共轭：IX.3.1 反变图式：IX.4.9 共变图式：IX.4.9 Coxeter 元：IX.4, 习题 14 支配：IX.7.1 特征值：VII.1.1 特征向量：VII.1.1

初等（自同构）：VII.3.1 嵌入（流形到子集邻域中的）：IX.9.1 例外型（可分裂单李代数）：VIII.3.2 指数：VIII.8.1 面（与抛物子集相联系的）：VIII.3.4 Fitting（模的分解）：VII.1.1 旗：VIII.13.1 Fourier 余变换：IX.8.1 Fourier 变换：IX.8.1 标架：IX.4.10 标架（分裂半单李代数的）：VIII.4.1 带标架的（半单李代数）：VIII.4.1 基本（支配权）：IX.7.1 基本（表示、单 g-模）：VIII.7.2 生成族（由标架定义的）：VIII.4.1 Harish-Chandra（同态）：VIII.6.4 Hermann Weyl（-- 的公式）：VIII.9.1 齐性对称空间：IX.1, 习题 8 不变（多项式函数）：VIII.8.3 对合（李代数）：IX.1, 习题 7 不可约（对合李代数）：IX.1, 习题 7 不可约（齐性对称空间）：IX.1, 习题 8 迷向（旗）：VIII.13.2, VIII.13.3 同构型（最高权 λ 的分支）：VIII.7.2 Jacobson-Morozov（定理）：VIII.11.2 格（Q-向量空间中的）：VIII.12.1 线性管：IX.9.3 极大环面：IX.2.2 极小（权）：VIII.7.3 缓增：IX.8.2 幂零（元素的分量）：VII.1.3 幂零空间：VII.1.1 结点（环面的群）：IX.4.2 结点向量：IX.4.5 序（Q-代数的）：VIII.12.1 正交（李代数）：VIII.13.2 正交（不可约表示）：VIII.7.5 抛物（子代数）：VIII.3.4, VIII.3.5 （正根的）分拆：VIII.9.1 许可（格）：VIII.12.6 正根：IX.7.1 准素（子空间）：VII.1.1 本原（模的元）：VIII.1.2, VIII,6.1 主（幂零元、单元素、sl₂ 三元组）：VIII.11.4 主（轨道）：IX.9.4 主（子群）：IX.4, 习题 18 真（子空间）：VII.1.1 拟极大（迷向旗）：VIII.13.4 R-极值（元）：VIII.7.2 R-饱和（子集）：VIII.7.2 根（子群）：IX.4.7 秩（李代数的）：VII.2.2 秩（李群的）：IX.2.4 速降：IX.8.2 实形式（复李代数的）：IX.3.1 既约（根图式）：IX.4.8 正则（线性表示的元）：VII.4.1 正则（李代数的元）：VII.2.2 正则（李群的元）：VII.4.2 表示环（李代数的）：VIII.7.6 根：IX.4.4 根（(g,h) 的）：VIII.2.2 根图式：IX.4.8 半单（元素的分量）：VII.1.3 半旋量（表示）：VIII.13.4 单（元素）：VIII.11.3 单（根）：IX.7.1 奇异超平面：IX.5.2 sl₂ 三元组：VIII.11.1 旋量（表示）：VIII.13.2, VIII.13.4 分裂（约化李代数）：VIII.2.1 可分裂（李子代数）：VII.5.1, VIII.10 可分裂（嘉当子代数、约化李代数）：VIII.2.1 可分裂包络（gl(V) 中子代数的）：VII.5.2 子群（嘉当子群）：IX.2.2 子群（极大秩子群）：IX.2.4 子群（主子群）：IX.4, 习题 18 辛（李代数）：VIII.13.3 辛（不可约表示）：VIII.7.5 切（多项式映射的线性映射）：VII.App.I.2 紧 Cʳ 收敛拓扑：IX.6.4 挠素数：IX.5, 习题 9 至 11 环面：IX.1.2 横截集：IX.9.3 及 IX.9，习题 16 型（An,Bn,\ldots）：IX.3.3 型（表示的）：IX.App.II.1, IX.App.II.2 向量（群）：IX.1.2 强正则（元）：IX.5.1

权（g-模的）：VII.1.1, VIII.1.2, VIII.6.1

权（表示的）：IX.4.3

权（分裂李代数的权群）：VIII.2.2

Weyl（群）：IX.2.5

Weyl（分裂代数的群）：VIII.2.2

Witt（基）：VIII.13.2, VIII.13.3, VIII.13.4

Zariski（拓扑）：VII.App.I.1